% Generate the table with make upload-data-verification (or CACHED=1).
\clearpage
\subsection{Data Verification}\label{app:data-verification}
\subsubsection{Where is the prescription upload data coming from?}
\label{subsec:upload-data-sources}
The report's upload analyses use
\texttt{EDLDB.BT\_HCA\_HCDM.rxp\_upload\_entity}.
Table~\ref{tab:upload-source-resource-types} compares the entries in that
source with \texttt{rxp\_notification\_service.upload\_entity}, resolved in
the \texttt{EDLDB} database, by \texttt{associated\_resource\_type}.

This verification covers all rows currently available in each source,
without date, CMPU, vet-diet, upload-type, or study-population filters.
Each physical row counts once; no deduplication or joins are applied.
These counts therefore describe the source tables rather than the report's
study cohort. They do not establish record-level equivalence between the
two sources or document their upstream lineage.

\input{../autogenerated/upload-data-verification/associated-resource-types.tex}

\noindent\textit{Notes:} Rows include every resource type present in either
source. A zero means that the source has no entries for that type.
\texttt{[SQL NULL]} denotes a missing resource type, separately from any
literal string value. Totals are queried independently from each source
and reconciled against the sum of its grouped counts before publication.
The extraction time above records the snapshot used for this comparison;
rerunning the query may change the counts as sources are updated.
\clearpage
